\documentclass[11pt]{article}
\usepackage[margin=1in]{geometry}
\usepackage{amsmath,amssymb,amsthm,hyperref}
\newtheorem{theorem}{Theorem}
\setlength{\emergencystretch}{3em}
\title{A symplectic-plane counterexample to the quadratic-difference equivalence}
\author{Myeongjin Kang}
\date{October 26, 2026}
\begin{document}\maketitle
\begin{abstract}
We give a complete Lean verification of a real two-dimensional counterexample to the quadratic-difference equivalence recorded as MathDB 358595. The zero endomorphism is a difference of two roots of $X^2+1$, but cannot be a difference of two roots that are alternating for a symplectic form. This is a concrete specialization of the counterexample already given in the cited primary manuscript, not a new resolution of an open conjecture.
\end{abstract}
\noindent\textit{Credit and license.} This AI-assisted note verifies the existing counterexample credited above. Original prose is CC BY 4.0 and original code is MIT.

An endomorphism $M$ is $b$-alternating if $b(x,Mx)=0$ for every vector $x$. A symplectic pair $(b,u)$ consists of a nondegenerate alternating bilinear form $b$ and a $b$-alternating endomorphism $u$. For monic quadratic polynomials $p,q$, an ordinary $(p,q)$-difference is an endomorphism $u=M-N$ with $p(M)=q(N)=0$. A symplectic $(p,q)$-difference further requires $M,N$ to be $b$-alternating. The conjectured equivalence asserts that, for every symplectic pair, the ordinary difference condition suffices for the symplectic one.
\begin{theorem}
Let $V=\mathbb R^2$, let $b(x,y)=x_0y_1-x_1y_0$, and let $p=q=X^2+1$. Then $(b,0)$ is a symplectic pair and $0$ is an ordinary $(p,p)$-difference, whereas $(b,0)$ is not a symplectic $(p,p)$-difference.
\end{theorem}
\begin{proof}
The expression for $b$ is bilinear and $b(x,x)=0$. If $b(x,y)=0$ for every $y$, taking $y=(0,1)$ and $y=(1,0)$ gives $x_0=x_1=0$; thus $b$ is nondegenerate. The zero map is plainly $b$-alternating. Both source polynomials are monic of degree two.

Set
\[
 J=\begin{pmatrix}0&1\\-1&0\end{pmatrix}.
\]
Since $J^2=-I$, we have $p(J)=0$, and the identity $0=J-J$ proves the ordinary difference condition.

Now suppose $M=(m_{ij})$ is $b$-alternating. For $x=(a,c)$,
\[
 b(x,Mx)=m_{10}a^2+(m_{11}-m_{00})ac-m_{01}c^2.
\]
Using the three vectors $(1,0),(0,1),(1,1)$ gives
$m_{10}=m_{01}=0$ and $m_{11}=m_{00}$. Hence $M=\alpha I$ for a real scalar $\alpha$. It follows that
\[
 p(M)=(\alpha^2+1)I\ne0.
\]
There is consequently no $b$-alternating root of $p$. In particular, a decomposition of zero into a difference of two such roots cannot exist.
\end{proof}
The formal module \texttt{QuadraticDifference.lean} constructs $b$ as an actual bilinear form, proves its alternation and nondegeneracy, and proves the polynomial monicity and degree. Its final predicates are on actual linear endomorphisms of $V$, not abstract matrix labels. The standard matrix-to-endomorphism algebra equivalence transports polynomial evaluation; the theorem \texttt{root\_end\_iff} proves this bridge. Its surjectivity ensures that the obstruction applies to every possible endomorphism decomposition. The main theorem is \texttt{QuadraticDifferenceCounterexample.source\_endomorphism\_counterexample}.

The source manuscript explicitly disproves its Conjecture 1 immediately after stating it in Section 1.3: over a field with an irreducible quadratic, it uses dimension $2n$ with $n$ odd. Our choice is $\mathbb R$, $n=1$, and $X^2+1$, with an elementary direct proof replacing the general invariant-factor argument. The author's MathDB reply also points readers to the existing counterexample. Thus MathDB's introductory suggestion that the conjecture remains open is incorrect. The present outcome is a full verified counterexample to the stated equivalence, with no novelty, priority, or first-formalization claim. Source snapshots, exact identifiers, and hashes accompany the certificate.
\begin{thebibliography}{9}
\bibitem{dsp} C. de Seguins Pazzis, \emph{The quadratic sum problem for symplectic pairs}, arXiv:2305.19340v1, Sections 1.2--1.3, Definition 1.1 and Conjecture 1 with its disproof. \url{https://arxiv.org/abs/2305.19340v1}.
\bibitem{mathdb} MathDB problem 358595 and reply by C. de Seguins Pazzis. \url{https://mathdb.com/p/358595}.
\end{thebibliography}
\end{document}
